% Options for packages loaded elsewhere
\PassOptionsToPackage{unicode}{hyperref}
\PassOptionsToPackage{hyphens}{url}
\documentclass[
  ignorenonframetext,
]{beamer}
\newif\ifbibliography
\usepackage{pgfpages}
\setbeamertemplate{caption}[numbered]
\setbeamertemplate{caption label separator}{: }
\setbeamercolor{caption name}{fg=normal text.fg}
\beamertemplatenavigationsymbolsempty
% remove section numbering
\setbeamertemplate{part page}{
  \centering
  \begin{beamercolorbox}[sep=16pt,center]{part title}
    \usebeamerfont{part title}\insertpart\par
  \end{beamercolorbox}
}
\setbeamertemplate{section page}{
  \centering
  \begin{beamercolorbox}[sep=12pt,center]{section title}
    \usebeamerfont{section title}\insertsection\par
  \end{beamercolorbox}
}
\setbeamertemplate{subsection page}{
  \centering
  \begin{beamercolorbox}[sep=8pt,center]{subsection title}
    \usebeamerfont{subsection title}\insertsubsection\par
  \end{beamercolorbox}
}
% Prevent slide breaks in the middle of a paragraph
\widowpenalties 1 10000
\raggedbottom
\AtBeginPart{
  \frame{\partpage}
}
\AtBeginSection{
  \ifbibliography
  \else
    \frame{\sectionpage}
  \fi
}
\AtBeginSubsection{
  \frame{\subsectionpage}
}
\usepackage{iftex}
\ifPDFTeX
  \usepackage[T1]{fontenc}
  \usepackage[utf8]{inputenc}
  \usepackage{textcomp} % provide euro and other symbols
\else % if luatex or xetex
  \usepackage{unicode-math} % this also loads fontspec
  \defaultfontfeatures{Scale=MatchLowercase}
  \defaultfontfeatures[\rmfamily]{Ligatures=TeX,Scale=1}
\fi
\usepackage{lmodern}
\ifPDFTeX\else
  % xetex/luatex font selection
\fi
% Use upquote if available, for straight quotes in verbatim environments
\IfFileExists{upquote.sty}{\usepackage{upquote}}{}
\IfFileExists{microtype.sty}{% use microtype if available
  \usepackage[]{microtype}
  \UseMicrotypeSet[protrusion]{basicmath} % disable protrusion for tt fonts
}{}
\makeatletter
\@ifundefined{KOMAClassName}{% if non-KOMA class
  \IfFileExists{parskip.sty}{%
    \usepackage{parskip}
  }{% else
    \setlength{\parindent}{0pt}
    \setlength{\parskip}{6pt plus 2pt minus 1pt}}
}{% if KOMA class
  \KOMAoptions{parskip=half}}
\makeatother
\usepackage{graphicx}
\makeatletter
\newsavebox\pandoc@box
\newcommand*\pandocbounded[1]{% scales image to fit in text height/width
  \sbox\pandoc@box{#1}%
  \Gscale@div\@tempa{\textheight}{\dimexpr\ht\pandoc@box+\dp\pandoc@box\relax}%
  \Gscale@div\@tempb{\linewidth}{\wd\pandoc@box}%
  \ifdim\@tempb\p@<\@tempa\p@\let\@tempa\@tempb\fi% select the smaller of both
  \ifdim\@tempa\p@<\p@\scalebox{\@tempa}{\usebox\pandoc@box}%
  \else\usebox{\pandoc@box}%
  \fi%
}
% Set default figure placement to htbp
\def\fps@figure{htbp}
\makeatother
\setlength{\emergencystretch}{3em} % prevent overfull lines
\providecommand{\tightlist}{%
  \setlength{\itemsep}{0pt}\setlength{\parskip}{0pt}}
% Slide layout defaults for warning-clean scientific decks.
\usepackage{etoolbox}
\IfFileExists{xurl.sty}{\usepackage{xurl}}{}
\IfFileExists{seqsplit.sty}{\usepackage{seqsplit}}{\newcommand{\seqsplit}[1]{#1}}
\protected\def\breakseq#1{\seqsplit{#1}}
\protected\def\breaktt#1{\begingroup\ttfamily\seqsplit{#1}\endgroup}
\setlength{\emergencystretch}{6em}
\tolerance=5000
\hbadness=10000
\hfuzz=1pt
\setlength{\tabcolsep}{2pt}
\AtBeginEnvironment{longtable}{\tiny\renewcommand{\arraystretch}{0.86}\setlength{\tabcolsep}{1pt}}
\AtBeginEnvironment{tabular}{\tiny\renewcommand{\arraystretch}{0.86}\setlength{\tabcolsep}{1pt}}

% Natbib and cross-reference fallbacks — slides don't load natbib
% or cleveref, but combined-PDF manuscript prose may emit these
% commands. The fallback renders citations as a bracketed key list
% and cross-references as detokenized labels so slides don't fail on
% undefined control sequences or raw underscores. \providecommand is
% a no-op if packages load later.
\providecommand{\citep}[1]{[#1]}
\providecommand{\citet}[1]{#1}
\providecommand{\citealp}[1]{#1}
\providecommand{\citeauthor}[1]{#1}
\providecommand{\citeyear}[1]{#1}
\providecommand{\cref}[1]{\texttt{\detokenize{#1}}}
\providecommand{\Cref}[1]{\texttt{\detokenize{#1}}}

\newtheorem{proposition}{Proposition}
\newtheorem{hypothesis}{Hypothesis}
\usepackage{bookmark}
\IfFileExists{xurl.sty}{\usepackage{xurl}}{} % add URL line breaks if available
\urlstyle{same}
% fallback for those not using the hyperref driver hyperxmp:
\makeatletter
\@ifundefined{xmpquote}{\newcommand{\xmpquote}[1]{#1}}{}
\makeatother
\hypersetup{
  hidelinks,
  pdfcreator={LaTeX via pandoc}}

\author{\texorpdfstring{}{}}
\date{}

\begin{document}

\section{Methods and analysis plan}\label{sec:methods}

\begin{frame}[fragile,allowframebreaks]{Methods and analysis plan}
The method is the registered-report workflow itself, implemented as
tested code. It has four locked stages that complete before any result
is observed, followed by two interpretation stages that run only after
the deviation ledger is built. \breakseq{[@fig:analysis\_dag]} shows this
ordering and the lock boundary it protects.

\begin{figure}
\centering
\pandocbounded{\includegraphics[keepaspectratio,alt={Registered-report analysis workflow, rendered from analysis\_plan\_stages() in src/registered\_report/demo\_study.py. The four blue stages --- freeze registration, validate completeness, collect/simulate data, and execute the registered analysis --- lock before results are revealed (dashed boundary). The two purple stages --- compare to plan and build the deviation ledger, then interpret confirmatory claims --- run only afterwards. The locks\_before\_results flag on each stage is what enforces the boundary in code.}]{../figures/analysis_dag.png}}
\caption{Registered-report analysis workflow, rendered from
\protect\breaktt{analysis\_plan\_stages()} in
\protect\breaktt{src/registered\_report/demo\_study.py}. The four blue
stages --- freeze registration, validate completeness, collect/simulate
data, and execute the registered analysis --- lock before results are
revealed (dashed boundary). The two purple stages --- compare to plan
and build the deviation ledger, then interpret confirmatory claims ---
run only afterwards. The \protect\breaktt{locks\_before\_results} flag on
each stage is what enforces the boundary in
code.}\label{fig:analysis_dag}
\end{figure}
\end{frame}

\begin{frame}[fragile,allowframebreaks]{Frozen registration}
\protect\phantomsection\label{frozen-registration}
\breaktt{freeze\_registration} deep-copies the registration, strips any
pre-existing hash, and stamps a canonical SHA-256 hash over the
sorted-key JSON encoding. \breaktt{validate\_registration} then checks
that \texttt{title}, \texttt{version}, \texttt{hypotheses},
\texttt{outcomes}, \texttt{exclusion\_rules}, and
\texttt{analysis\_plan} are present; that each hypothesis has a unique
\texttt{id} and a \texttt{claim}; that each outcome names a
\texttt{measure} and \texttt{analysis}; that
\breaktt{analysis\_plan.primary\_model} and \texttt{seed} are set; and
that registered-report stage and ethics-review metadata are supplied.
For the demonstration registration this validation returns no findings.
\end{frame}

\begin{frame}[fragile,allowframebreaks]{Registered analysis plan}
\protect\phantomsection\label{registered-analysis-plan}
The analysis plan is fixed before execution:

\begin{itemize}
\tightlist
\item
  \textbf{Primary model} --- a two-sided label-permutation test on the
  group mean difference.
\item
  \textbf{Significance threshold} --- \texttt{alpha\ =\ 0.05}.
\item
  \textbf{Seed} --- \texttt{20260709}, used for both data synthesis and
  the permutation schedule so the result is reproducible.
\item
  \textbf{Permutations} --- \texttt{2000} label shuffles; the two-sided
  p-value uses the add-one correction
  \breaktt{(k\ +\ 1)\ /\ (permutations\ +\ 1)} and is therefore strictly
  positive.
\end{itemize}
\end{frame}

\begin{frame}[fragile,allowframebreaks]{Deterministic demonstration
data}
\protect\phantomsection\label{deterministic-demonstration-data}
Because this is a template, the ``data collection'' stage is a seeded
synthetic generator (\breaktt{generate\_demo\_dataset}), \textbf{not} a
real study. It draws \texttt{n\ =\ 24} control observations from
\texttt{Normal(0,\ 1)} and \texttt{n\ =\ 24} treatment observations from
\texttt{Normal(0.8,\ 1)} using a single seeded generator, so the entire
dataset is a pure function of the plan seed.
\breaktt{run\_registered\_analysis} reads the seed and \texttt{alpha}
back out of the frozen plan and executes the registered permutation test
against this dataset, guaranteeing the executed analysis honours the
locked plan. The executed summary is written to
\breaktt{output/data/demo\_analysis.json} and is the single source of
truth for every number in {[}Section @sec:results{]}.

Forks replace this generator with genuine data collection while keeping
the same \breaktt{run\_registered\_analysis} contract, so the
confirmatory analysis remains bound to the frozen plan rather than to
post-hoc choices.
\end{frame}

\end{document}
